% ============================================================
% 第 12 章 三相感应电机概述（教材第四篇，印刷页 125--130）
% 节名：第一节 三相感应电机的运行状态；第二节 三相感应电机的结构；第三节 三相感应电动机的额定值
% ============================================================
\chapter{三相感应电机概述}

感应电机（异步电机）多用作电动机，在少数场合也作发电机用。它是产量最大、应用最广的一种电机。本章只讲"是什么、怎么转、铭牌怎么看"，为第 13、14 章的理论分析做准备。

\section{三相感应电机的运行状态}

感应电机的定子接三相电源产生旋转磁场，转子电路自成闭合回路（笼型）或经外接电阻闭合（绕线型）。转子转速 $n$ 与旋转磁场同步转速 $n_1$ 之间的差异，决定了电机的运行状态。定义\textbf{转差率}
\begin{equation}
s=\frac{n_1-n}{n_1},
\end{equation}
则 $n=(1-s)n_1$。

\keyformula{s=\frac{n_1-n}{n_1},\qquad n=(1-s)n_1}

按转差率的正负与大小，感应电机有三种运行状态。

\subsection{电动机状态（$0<s<1$）}
$0<n<n_1$。转子切割旋转磁场的方向与磁场转向相同，转子感应电动势与电流产生的电磁转矩方向与旋转磁场转向相同，是\textbf{驱动}转矩；电机从电网吸收电功率，扣除损耗后以机械功率输出。这是感应电机最主要的运行状态。

\subsection{发电机状态（$s<0$）}
$n>n_1$。转子被原动机拖动超过同步转速，转子导体相对旋转磁场反向切割，感应电动势、电流反向，电磁转矩变为\textbf{制动}转矩，即吸收机械功率、输出电功率。感应发电机需要电网提供励磁（吸收无功），不能单独建压，故应用不如同步发电机广泛，但在风力发电中大量使用（双馈、笼型异步）。

\subsection{电磁制动状态（$s>1$）}
$n<0$（转子逆磁场方向旋转）。此时转子导体切割磁场的方向虽仍使电流产生制动转矩，但电机既吸收电网电功率、又吸收机械功率，两者都变成转子回路的损耗热——这是"反接制动"状态，用于快速停车，不能长期运行。

\begin{center}
\begin{tabular}{@{}llll@{}}
\toprule
运行状态 & 转速范围 & 转差率 & 功率流向 \\
\midrule
电动机 & $0<n<n_1$ & $0<s<1$ & 电网 $\rightarrow$ 机械（输出） \\
发电机 & $n>n_1$ & $s<0$ & 机械 $\rightarrow$ 电网（输出） \\
电磁制动 & $n<0$ & $s>1$ & 电网 $+$ 机械 $\rightarrow$ 转子损耗（发热） \\
\bottomrule
\end{tabular}
\end{center}

\corepoint{转差率是感应电机唯一的"状态变量"：$s$ 一变，转子频率、转子电动势、电流、转矩、功率分配全变。牢记三句话：$s=0$（同步转速）时转子无感应电动势、无转矩；$s=1$（堵转/起动瞬间）时转子电流最大、起动转矩有限；$s$ 略大于 0（额定运行）时效率最高。凡是"转速变化引起什么变化"的问题，一律先写出 $s$，再由 $s$ 推其余量。}

\section{三相感应电机的结构}

感应电机由\textbf{定子}、\textbf{转子}和定转子之间的\textbf{气隙}三部分组成。

\subsection{定子}
定子由定子铁芯、定子绕组和机座构成。
\begin{itemize}
  \item \textbf{定子铁芯}：由 0.5\,mm 厚的硅钢片叠成，内圆均匀开槽，用以嵌放定子绕组，同时构成主磁路。
  \item \textbf{定子绕组}：三相对称绕组，可接成星形或三角形，其六个出线端引至接线盒，便于改变接法（这是 Y—$\Delta$ 降压起动的物理基础）。
  \item \textbf{机座}：铸铁或钢板焊接，起支撑与防护作用；封闭式电机机座外表面带散热筋。
\end{itemize}

\subsection{转子}
按转子绕组形式分为\textbf{笼型}和\textbf{绕线型}两种。

\subsubsection{笼型转子}
转子铁芯外圆开槽，槽内嵌放\textbf{导条}（铜条或铸铝），导条两端用\textbf{端环}短接，形成自行闭合的多相绕组，形似鼠笼。笼型转子结构简单、坚固耐用、价格低廉、无需维护，是应用最广的形式；缺点是起动转矩较小、起动电流较大，且转子回路无法外接电阻。

\subsubsection{绕线型转子}
转子铁芯槽内嵌放三相对称绕组，一般接成星形，三个出线端通过\textbf{滑环}和电刷引至外电路，可串入外加电阻。绕线型转子的优点是可改善起动性能（串电阻提高起动转矩、减小起动电流）并可实现调速（串电阻调速、串级调速）；缺点是结构复杂、成本高、需维护滑环电刷。

\subsection{气隙}
感应电机的气隙很小（中小型电机一般为 0.2～1.5\,mm），因为气隙越大，励磁电流越大、功率因数越低。这是感应电机与同步电机结构上的一个重要差别（同步电机气隙可较大）。

\begin{center}
\begin{tabular}{@{}p{2.4cm}p{4.6cm}p{4.6cm}@{}}
\toprule
项目 & 笼型转子 & 绕线型转子 \\
\midrule
转子绕组 & 导条 + 端环，自行闭合 & 三相绕组，经滑环引出 \\
外接电阻 & 不能 & 可以 \\
起动性能 & 起动电流大、起动转矩小 & 串电阻可提高 $T_{st}$、减小 $I_{st}$ \\
调速 & 主要靠变频、变极 & 可转子串电阻、串级调速 \\
结构与维护 & 简单、坚固、免维护 & 复杂，滑环电刷需维护 \\
典型用途 & 风机、水泵、机床（绝大多数） & 起重、轧机、需大起动转矩场合 \\
\bottomrule
\end{tabular}
\end{center}

\section{三相感应电动机的额定值}

感应电动机铭牌上标明额定值和一些与运行有关的技术数据。按铭牌指定的条件和额定值运行，就叫额定运行。

\subsection{额定功率 $P_N$}
额定运行时电动机轴上输出的\textbf{机械功率}，单位 kW。注意这是\textbf{输出}功率（轴功率），不是输入电功率。

\subsection{额定电压 $U_N$}
额定运行时定子绕组应加的\textbf{线电压}，单位 V 或 kV。

\subsection{额定电流 $I_N$}
额定运行时定子绕组的\textbf{线电流}，单位 A。

\subsection{额定转速 $n_N$}
额定运行时的转子转速，单位 r/min。由 $n_N$ 可算出额定转差率
\begin{equation}
s_N=\frac{n_1-n_N}{n_1},
\end{equation}
其中同步转速 $n_1=60f_N/p$ 由额定频率和极对数确定。

\subsection{额定频率 $f_N$}
我国为 50\,Hz。

\subsection{额定功率因数 $\cos\varphi_N$}
额定运行时定子侧的功率因数。

\subsection{额定效率 $\eta_N$}
额定运行时输出功率与输入功率之比，故输入功率
\begin{equation}
P_{1N}=\frac{P_N}{\eta_N}=\sqrt3U_NI_N\cos\varphi_N .
\end{equation}

\keyformula{P_{1N}=\sqrt3U_NI_N\cos\varphi_N=\frac{P_N}{\eta_N}}

\subsection{接法与绝缘等级}
铭牌还标明定子绕组的接法（Y 或 $\Delta$）与绝缘等级。接法必须与电源线电压相匹配：例如铭牌"380\,V/$\Delta$ 接"，表示电源线电压 380\,V 时应接成三角形。

\begin{center}
\begin{tabular}{@{}lll@{}}
\toprule
额定值 & 含义 & 注意 \\
\midrule
$P_N$ & 轴上输出的机械功率 & 输出，不是输入 \\
$U_N$ & 定子线电压 & 线值 \\
$I_N$ & 定子线电流 & 线值 \\
$n_N$ & 额定转速 & 略低于 $n_1$，$s_N$ 约 1\%～5\% \\
$f_N$ & 额定频率 & 我国 50\,Hz \\
$\cos\varphi_N$ & 定子功率因数 & 中小型约 0.8～0.9 \\
$\eta_N$ & 额定效率 & 中小型约 85\%～93\% \\
接法 & Y 或 $\Delta$ & 必须与电源电压匹配 \\
\bottomrule
\end{tabular}
\end{center}

\corepoint{额定值计算最容易错的两处：① 把 $P_N$ 当成输入功率——由 $P_N$、$\eta_N$ 求 $P_1$ 时才乘/除效率，方向别反；② 用 $U_N$、$I_N$ 算功率时别忘了 $\sqrt3$（线值），而在等效电路里算每相损耗时又要把线值换成相值。记住一条链：$P_1=\sqrt3U_NI_N\cos\varphi_N$，$P_1=P_N+p_{Cu1}+p_{Fe}+p_{Cu2}+p_{mec}+p_{ad}$，$P_N=P_1-\sum p$。}
